\documentclass[../main.tex]{subfiles}
\begin{document}

Chương này giới thiệu tổng quan về đề tài. Trước hết, chương đặt vấn đề thông qua bối cảnh
hoạt động điểm danh tại các trường đại học cùng những hạn chế của các phương pháp hiện hành.
Tiếp đó, chương xác định mục tiêu và phạm vi của đề tài, trình bày định hướng giải pháp cùng
các đóng góp chính, và cuối cùng mô tả bố cục của toàn bộ quyển báo cáo.

% ─────────────────────────────────────────────
\section{Đặt vấn đề}
\label{sec:dat-van-de}

Điểm danh là một hoạt động bắt buộc và diễn ra thường xuyên trong công tác quản lý đào
tạo tại các trường đại học. Kết quả điểm danh không chỉ phản ánh ý thức học tập của
sinh viên mà còn là một căn cứ để đánh giá điều kiện dự thi theo quy chế đào tạo. Tuy
nhiên, cho đến nay phần lớn các lớp học vẫn sử dụng những phương pháp điểm danh truyền
thống như gọi tên, ký tên vào danh sách giấy, hoặc chuyền tay phiếu điểm danh. Các
phương pháp này tuy đơn giản nhưng tiêu tốn nhiều thời gian của giờ học, đặc biệt với
các lớp học phần đông sinh viên, và rất khó kiểm soát tính trung thực.

Nhằm khắc phục các nhược điểm trên, nhiều giải pháp điểm danh ứng dụng công nghệ đã
được nghiên cứu và triển khai, bao gồm thẻ từ và thẻ RFID, công nghệ NFC, đèn hiệu
Bluetooth năng lượng thấp (BLE Beacon), mã QR, hay nhận diện khuôn mặt. Mỗi giải pháp
giải quyết được một phần bài toán nhưng đồng thời cũng bộc lộ những hạn chế riêng về
chi phí phần cứng, độ phức tạp khi triển khai hoặc khả năng chống gian lận, sẽ được
phân tích chi tiết trong Chương~\ref{chapter:Related_works}.

Song song với sự phát triển của các giải pháp công nghệ, sự phổ biến của điện thoại
thông minh và các nền tảng nhắn tin đã tạo ra một cơ hội mới. Tại Việt Nam, Zalo là một
trong những ứng dụng nhắn tin phổ biến nhất với hơn 78 triệu người dùng thường xuyên hằng
tháng~\cite{zalostats}, gần như bao phủ toàn bộ nhóm sinh viên. Nền tảng \textit{Zalo Mini App} cho
phép xây dựng các ứng dụng chạy trực tiếp bên trong Zalo mà người dùng không cần cài
đặt thêm phần mềm, đồng thời cung cấp sẵn cơ chế xác thực người dùng, truy cập máy ảnh,
định vị và quét mã QR~\cite{zalominiapp}. Đây là tiền đề thuận lợi để xây dựng một hệ
thống điểm danh hiện đại, tiếp cận được số đông sinh viên với chi phí triển khai thấp.

Từ thực tiễn nêu trên, đề tài tập trung giải quyết bốn vấn đề chính mà các hệ thống
điểm danh hiện hành tại trường đại học còn gặp phải:

\begin{itemize}[leftmargin=1.5em]
  \item \textbf{Gian lận điểm danh hộ (proxy attendance).} Đây là vấn đề nghiêm trọng
  nhất: một sinh viên có thể điểm danh thay cho bạn vắng mặt bằng cách ký hộ, chụp ảnh
  màn hình mã QR gửi cho nhau, hoặc khai báo hộ thông tin. Phần lớn các phương pháp hiện
  hành không có cơ chế xác minh danh tính người điểm danh tại thời điểm điểm danh.

  \item \textbf{Tốc độ chậm với lớp học đông.} Với các lớp học phần trên 50 sinh viên,
  việc gọi tên hoặc ký danh sách có thể chiếm nhiều phút đầu giờ, gây gián đoạn bài
  giảng và giảm hiệu quả sử dụng thời gian trên lớp.

  \item \textbf{Giảng viên khó giám sát theo thời gian thực.} Giảng viên không có công
  cụ để theo dõi tiến trình điểm danh ngay khi đang diễn ra, cũng như thiếu báo cáo
  thống kê dài hạn phục vụ việc đánh giá chuyên cần.

  \item \textbf{Thiếu kênh xin phép nghỉ chính thức.} Sinh viên không có cách gửi đơn
  xin nghỉ có lý do trong hệ thống, dẫn đến khó khăn khi đối chiếu và xử lý các trường
  hợp vắng có phép.
\end{itemize}

Những vấn đề này đặt ra câu hỏi nghiên cứu trọng tâm của đề tài: \textit{làm thế nào
kết hợp mã QR động, xác minh sinh trắc học và xác minh ngang hàng giữa các sinh viên để
xây dựng một quy trình điểm danh vừa nhanh chóng, thuận tiện, vừa có khả năng chống gian
lận cao mà không đòi hỏi phần cứng chuyên dụng đắt tiền?} Việc trả lời câu hỏi này chính
là động lực thúc đẩy việc thực hiện đề tài.

% ─────────────────────────────────────────────
\section{Mục tiêu và phạm vi đề tài}
\label{sec:muc-tieu}

\textbf{Mục tiêu tổng quát} của đồ án là xây dựng một hệ thống điểm danh thông minh,
chống gian lận, hoạt động trên nền tảng Zalo Mini App, áp dụng được cho môi trường đào
tạo đại học. Hệ thống có tên \textit{Zimo Check-in}.

Để đạt được mục tiêu tổng quát, đồ án đặt ra các \textbf{mục tiêu cụ thể} sau:

\begin{enumerate}[leftmargin=1.5em]
  \item Xây dựng ứng dụng điểm danh trên Zalo Mini App sử dụng mã QR động xoay theo chu
  kỳ 30 giây dựa trên thuật toán HMAC-SHA256 nhằm chống việc chụp màn hình và chia sẻ mã.

  \item Tích hợp cơ chế xác minh khuôn mặt phía thiết bị người dùng để xác thực danh
  tính sinh viên tại thời điểm điểm danh.

  \item Thiết kế cơ chế xác minh ngang hàng (peer-to-peer) yêu cầu mỗi sinh viên trao
  đổi mã xác nhận với tối thiểu ba sinh viên xung quanh, nhằm chống điểm danh từ xa.

  \item Xây dựng trang quản trị web cho giảng viên và phòng đào tạo phục vụ quản lý lớp,
  giám sát phiên điểm danh theo thời gian thực và xuất báo cáo.

  \item Thiết kế hệ thống phát hiện gian lận nhiều lớp với khả năng cảnh báo và lập báo
  cáo tự động.
\end{enumerate}

\textbf{Về nghiệp vụ}, hệ thống phục vụ ba nhóm vai trò người dùng: sinh viên, giảng
viên và quản trị viên (đại diện cho phòng đào tạo).

\textbf{Về kỹ thuật}, phạm vi triển khai gồm: ứng dụng phía người dùng là một Zalo Mini
App xây dựng bằng React, TypeScript và thư viện quản lý trạng thái Jotai; trang quản trị
là một ứng dụng web; và phần xử lý phía máy chủ dựa trên nền tảng Firebase (Firestore,
Cloud Functions, Storage).

\textbf{Về quy mô triển khai}, đồ án được thực hiện ở mức thử nghiệm (pilot) trong khuôn
khổ một đồ án tốt nghiệp, chưa triển khai chính thức ở quy mô toàn trường.

\textbf{Các nội dung nằm ngoài phạm vi} của đồ án bao gồm: tích hợp với hệ thống quản lý
học tập (LMS) của nhà trường, tích hợp thẻ sinh viên RFID, và cơ chế chống giả mạo khuôn
mặt bằng kỹ thuật phát hiện sự sống (liveness detection) ở mức nâng cao. Những hạn chế
này được trình bày trung thực ở Chương~\ref{chapter:conclusion}.

% ─────────────────────────────────────────────
\section{Định hướng giải pháp}
\label{sec:dong-gop}

Để giải quyết những vấn đề nêu trên, đề tài lựa chọn định hướng xây dựng hệ thống điểm danh
ngay trên nền tảng Zalo Mini App, tận dụng việc sinh viên đã sẵn có Zalo nên không cần cài
đặt ứng dụng mới. Thay vì dựa vào một yếu tố xác minh duy nhất, hệ thống kết hợp nhiều lớp
xác minh bổ trợ (mã QR động, khuôn mặt và ngang hàng) theo nguyên lý phòng thủ theo chiều
sâu; các nghiệp vụ nhạy cảm được tập trung xử lý phía máy chủ trên nền tảng Firebase nhằm
bảo đảm an toàn. Theo định hướng đó, các đóng góp chính của đồ án bao gồm:

\begin{enumerate}[leftmargin=1.5em]
  \item \textbf{Cơ chế mã QR động dựa trên HMAC-SHA256 xoay 30 giây}, có đồng bộ đồng
  hồ giữa nhiều thiết bị (điện thoại của giảng viên và màn hình máy chiếu), giúp chống
  hiệu quả việc chụp và chia sẻ lại mã QR.

  \item \textbf{Tính năng ``Cổng máy chiếu'' (Teacher Display)} cho phép hiển thị mã QR
  cỡ lớn trên máy chiếu của lớp học, sử dụng cơ chế ghép nối (pairing) bằng mã thông báo
  128-bit với quy tắc chuyển trạng thái một chiều, an toàn.

  \item \textbf{Quy trình điểm danh đa yếu tố bốn bước} (mã QR, khuôn mặt, ngang hàng,
  định vị) kết hợp thành một hệ thống chống gian lận nhiều lớp theo nguyên lý phòng thủ
  theo chiều sâu (defense in depth).

  \item \textbf{Kiến trúc ba vai trò với phân quyền chặt chẽ}, sử dụng quy tắc bảo mật
  của Firestore kết hợp cơ chế cấp token tùy biến (Custom Token).

  \item \textbf{Bảo đảm tuân thủ chính sách của nền tảng Zalo}, bao gồm chính sách quyền
  riêng tư, xử lý dữ liệu sinh trắc học và cơ chế đồng thuận của người dùng.
\end{enumerate}

% ─────────────────────────────────────────────
\section{Bố cục đồ án}
\label{sec:bo-cuc}

Phần còn lại của quyển báo cáo được tổ chức như sau:

\begin{itemize}[leftmargin=1.5em]
  \item \textbf{Chương~\ref{chapter:Related_works} -- Khảo sát và phân tích yêu cầu:}
  khảo sát các hệ thống điểm danh hiện có, so sánh ưu nhược điểm, phân tích đối tượng
  người dùng và xác định các yêu cầu chức năng, phi chức năng.

  \item \textbf{Chương~\ref{chapter:Methodology} -- Nền tảng lý thuyết và công nghệ sử
  dụng:} trình bày các công nghệ và kiến thức nền tảng được áp dụng, gồm nền tảng Zalo
  Mini App, hệ sinh thái Firebase, thuật toán HMAC-SHA256, nhận diện khuôn mặt và mã QR.

  \item \textbf{Chương~\ref{chapter:Experiment} -- Phân tích thiết kế, triển khai và
  đánh giá hệ thống:} trình bày kiến trúc tổng thể, thiết kế cơ sở dữ liệu, các thuật
  toán cốt lõi, quá trình triển khai và kết quả đánh giá về chức năng, bảo mật, hiệu năng.

  \item \textbf{Chương~\ref{chapter:SolutionAndContribution} -- Các giải pháp và đóng
  góp nổi bật:} phân tích sâu những giải pháp kỹ thuật tạo nên sự khác biệt của hệ thống.

  \item \textbf{Chương~\ref{chapter:conclusion} -- Kết luận và hướng phát triển:} đánh
  giá lại những gì đã làm được, chỉ ra các điểm còn hạn chế và phác thảo lộ trình mở rộng.
\end{itemize}

% ─────────────────────────────────────────────
\subsection*{Kết chương}
Chương này đã đặt vấn đề từ bối cảnh và các khó khăn thực tiễn của hoạt động điểm danh, xác
định mục tiêu, phạm vi, trình bày định hướng giải pháp cùng những đóng góp chính của đề tài,
đồng thời giới thiệu bố cục của quyển báo cáo. Chương tiếp theo sẽ khảo sát các hệ thống điểm
danh hiện có và phân tích các yêu cầu đặt ra cho hệ thống cần xây dựng.

\end{document}
